\chapter{입력은 몇 개나 필요한가}
% full version: chapters/19_dendrites.tex

뉴런에는 가지를 친 입력 돌기, 곧 가지돌기가 있다. 어떤 뉴런에서는 접점이 수천 개인 나무 한 그루이고, 어떤 뉴런에서는 가지 한두 개이거나 아예 없다. 엔지니어에게 입력의 수는 회로의 파라미터이며, 과제에 맞춰 골라진 값이다.

\section*{0개에서 10만 개까지}

촉각 감지기와 통증 감지기에는 가지돌기가 전혀 없다. 전선이 피부에서 척수까지 그냥 뻗어 있고, 감지기는 그 끝에 붙어 있다. 먼 끝에 감지기가 달린 선로다.

\pencilfig{19}{\pencilart{19}}{입력이 많으면 가산기, 거대한 입력 하나이면 빠르고 정확한 중계기다.}

후각 뉴런, 그리고 눈과 귀의 감지기에서 신호를 넘겨받는 뉴런에는 가지돌기가 하나다. 입력 하나, 출력 하나. 리피터이자 버퍼다.

소리 방향 검출기에는 서로 반대쪽을 향한 가지돌기가 둘 있다. 귀마다 하나씩이다. 두 입력을 서로 다른 가지에 나누어 받으면 뉴런의 “그리고”가 더 엄격해진다. 한 가지에 두 전류가 함께 들어오면 서로 방해할 것이기 때문이다.

움직임 학습 회로의 작은 뉴런에는 짧은 가지돌기가 네 개쯤 있고, 저마다 입력 하나를 붙잡는다. 이런 뉴런이 수백억 개이며, 각각은 신호를 네 개만 더한다. 그 대신 이들은 함께 입력을 수백만 가지 조합으로 펼치고, 입력이 10만 개인 큰 출력 뉴런이 그 가운데 필요한 것을 고른다.

마지막으로 거대한 나무와 수천 개의 입력을 가진 뉴런은 가산기이자 분류기다. 이런 뉴런에서는 가지 하나하나가 거의 독립된 작은 회로처럼 일한다.

\section*{왜 그런가}

모든 것은 충전이 결정한다. 입력이 적은 작은 뉴런은 큰 연결 하나로 1밀리초도 안 되어 충전된다. 이런 뉴런은 타이밍이 정확하고 신호를 믿을 만하게 넘겨준다. 큰 나무는 연결 하나로는 결코 충전되지 않는다. 그 전에 전하가 새어 나간다. 큰 나무에는 함께 도착한 신호 수십 개가 필요하지만, 그 대신 신호들을 더하고 비교할 수 있다. 입력이 적고 연결이 크면 정확한 타이밍에 쓰인다. 입력이 많으면 셈에 쓰인다.

\pencilfig{128}{%
% charging: a small neuron reaches threshold from one large synapse at once; a big tree barely moves
% from one input and needs many inputs arriving together
\foreach \dx/\t in {0/{입력 적음}, 4.3/{큰 나무}}{
  \draw[pen] (\dx,0) -- (\dx+3.6,0);
  \draw[pcGray, densely dashed, line width=0.5pt] (\dx,0.9) -- (\dx+3.6,0.9);
  \node[font=\small\itshape, text=pcGray, anchor=south] at (\dx+1.8,1.75) {\t};}
\node[lbl, anchor=east] at (-0.05,0.9) {문턱값};
% small neuron
\draw[pcBlue, line width=2pt] (0.6,-0.15) -- (0.6,-0.6);
\draw[penlite=pcBlue] (0,0) -- (0.6,0) -- plot[domain=0.6:0.8, samples=12] (\x,{1.3*(1-exp(-(\x-0.6)/0.18))}) -- (0.8,1.6) -- (0.84,0) -- (3.5,0);
\node[lbl, text=pcRed, anchor=west] at (0.85,1.45) {곧바로 딸깍};
\node[lbl, text=pcBlue, anchor=north] at (0.6,-0.62) {큰 입력 하나};
% big tree
\draw[pcBlue, line width=0.6pt] (4.8,-0.15) -- (4.8,-0.45);
\foreach \s in {6.15,6.2,6.25,6.3,6.35,6.4,6.45,6.5}{\draw[pcBlue, line width=0.6pt] (\s,-0.15) -- (\s,-0.45);}
\draw[penlite=pcBlue] (4.3,0) -- (4.8,0) -- plot[domain=4.8:6.15, samples=20] (\x,{0.16*((\x-4.8)/0.15)*exp(1-(\x-4.8)/0.15)}) -- plot[domain=6.15:6.62, samples=14] (\x,{1.25*(1-exp(-(\x-6.15)/0.4))}) -- (6.62,1.6) -- (6.66,0) -- (7.9,0);
\node[lbl, anchor=south] at (4.95,0.2) {거의 그대로};
\node[lbl, text=pcBlue, anchor=north] at (6.33,-0.47) {여럿이 함께};
}{작은 뉴런은 큰 연결 하나로 1밀리초도 안 되어 충전되고 곧바로 딸깍한다. 큰 나무는 입력 하나로는 거의 변하지 않는다. 함께 도착한 신호 수십 개가 필요하다.}

이렇게 해서 뉴런을 분류하는 세 번째 방법을 얻었다. 물질에 따라, 소자에 따라, 입력의 수에 따라. 각 분류는 같은 기계를 새로운 쪽에서 비춘다.

\fullversion{19}
